% Common mathematical setup for the research templates.
% Keep this file focused on notation that is useful across economics and
% mathematics documents; load domain-specific packages in the document header.

% Core AMS and research-math packages.
\usepackage{amsmath,amsfonts,amssymb,mathtools,mathrsfs}
\usepackage{bm}
\usepackage{bbm}
\usepackage{graphicx}

% Mathematical diagrams. This is retained because the Note template uses
% tikz-cd; it can be commented out for a text-only paper.
\usepackage{tikz-cd}
\usetikzlibrary{calc}
\usetikzlibrary{decorations.pathmorphing}

% A TikZ style for curved arrows of a fixed height, due to AndréC.
\tikzset{curve/.style={settings={#1},to path={(\tikztostart)
                        .. controls ($ (\tikztostart)!\pv{pos}!(\tikztotarget)!\pv{height}!270:(\tikztotarget) $)
                        and ($ (\tikztostart)!1-\pv{pos}!(\tikztotarget)!\pv{height}!270:(\tikztotarget) $)
                        .. (\tikztotarget)\tikztonodes}},
        settings/.code={\tikzset{quiver/.cd,#1}
                \def\pv##1{\pgfkeysvalueof{/tikz/quiver/##1}}},
        quiver/.cd,pos/.initial=0.35,height/.initial=0}

% TikZ arrowhead/tail styles.
\tikzset{tail reversed/.code={\pgfsetarrowsstart{tikzcd to}}}
\tikzset{2tail/.code={\pgfsetarrowsstart{Implies[reversed]}}}
\tikzset{2tail reversed/.code={\pgfsetarrowsstart{Implies}}}
\tikzset{no body/.style={/tikz/dash pattern=on 0 off 1mm}}

% Common notation.
\providecommand{\at}[3]{\left.#1\right\vert_{#2}^{#3}}
\providecommand\quotient[2]{%
    \mathchoice
    {% \displaystyle
        \text{\raise1ex\hbox{$#1$}\Big/\lower1ex\hbox{$#2$}}%
    }
    {% \textstyle
        #1\,/\,#2
    }
    {% \scriptstyle
        #1\,/\,#2
    }
    {% \scriptscriptstyle
        #1\,/\,#2
    }%
}

% Number systems and research notation.
\providecommand{\N}{\mathbb{N}}
\providecommand{\Z}{\mathbb{Z}}
\providecommand{\Q}{\mathbb{Q}}
\providecommand{\R}{\mathbb{R}}
\providecommand{\C}{\mathbb{C}}
\providecommand{\E}{\mathbb{E}}
\providecommand{\Prob}{\mathbb{P}}
\providecommand{\one}{\mathbbm{1}}
\providecommand{\indep}{\mathrel{\perp\!\!\!\perp}}
\providecommand{\iid}{\overset{\mathrm{iid}}{\sim}}
\providecommand{\given}{\middle|}
\providecommand{\abs}[1]{\left\lvert#1\right\rvert}
\providecommand{\norm}[1]{\left\lVert#1\right\rVert}
\providecommand{\set}[1]{\left\{#1\right\}}
\providecommand{\mat}[1]{\mathbf{#1}}
\providecommand{\vect}[1]{\bm{#1}}
\providecommand{\diff}{\mathrm{d}}
\providecommand{\deriv}[2]{\frac{\diff #1}{\diff #2}}
\providecommand{\pderiv}[2]{\frac{\partial #1}{\partial #2}}

\providecommand{\True}{\textsf{True}}
\providecommand{\T}{\textsf{T}}
\providecommand{\False}{\textsf{False}}
\providecommand{\F}{\textsf{F}}

% Operators used frequently in theoretical and empirical work.
\DeclareMathOperator{\id}{id}
\DeclareMathOperator{\Span}{span}
\DeclareMathOperator{\supp}{supp}
\DeclareMathOperator{\vol}{vol}
\DeclareMathOperator{\dist}{dist}
\DeclareMathOperator{\diam}{diam}
\DeclareMathOperator{\im}{Im}
\DeclareMathOperator{\sgn}{sgn}
\DeclareMathOperator{\Int}{Int}
\DeclareMathOperator{\diag}{diag}
\DeclareMathOperator{\dom}{dom}
\DeclareMathOperator{\plim}{plim}
\DeclareMathOperator{\rank}{rank}
\DeclareMathOperator{\tr}{tr}
\DeclareMathOperator*{\argmax}{arg\,max}
\DeclareMathOperator*{\argmin}{arg\,min}
\DeclareMathOperator{\Var}{Var}
\DeclareMathOperator{\Cov}{Cov}
\DeclareMathOperator{\Corr}{Corr}

% Do not globally redefine AMS relation commands such as \implies and \iff.
% The standard definitions provide the correct relation spacing.

% Optional specialized packages. Uncomment only when a document uses them.
% \usepackage{physics}   % broad derivative/vector shortcuts; may conflict with local notation
% \usepackage{complexity} % computational-complexity notation. Disabled here: it
% redefines \C, which already denotes the complex numbers above, and it plays no
% role in a mathematics-for-economists volume.
\usepackage{stmaryrd}  % extended semantic symbols, including \lightning
\newcommand\conta{\scalebox{1.1}{\(\lightning\)}}
\providecommand{\act}{\rotatebox[origin=c]{-180}{\(\,\circlearrowright\,)}}

% Figure support for Inkscape's PDF+LaTeX export.
\usepackage{import}
\usepackage{xifthen}
\pdfminorversion=7 % only needed when a downstream figure requires PDF 1.7
\usepackage{pdfpages}    % optional support for inserting complete PDF pages
\usepackage{transparent} % optional transparency support for imported graphics
\providecommand{\incfig}[2][\columnwidth]{%
    \def\svgwidth{#1}%
    \import{Figures/}{#2.pdf_tex}%
}
